% Bench layout: one shield, one marker wire, one question to settle in UM3120.
\begin{tikzpicture}[font=\sffamily\scriptsize]
\node[board, minimum width=52mm, minimum height=28mm] (nuc) at (0,0) {};
\node[text=white, font=\sffamily\bfseries\small] at (-0.55,-1.05) {NUCLEO-H7A3ZI-Q};

\node[pinrow, minimum width=42mm] at (0,1.25) {};
\node[pinrow, minimum width=42mm] at (0,-0.55) {};
\foreach \x in {-2.05,-1.85,...,2.05}{\fill[yellow!80!black] (\x,1.25) circle (0.35mm);}

\node[hat, minimum width=38mm, minimum height=17mm] (sh) at (0.1,0.45) {};
\node[text=white, font=\sffamily\bfseries\small] at (0.1,0.95) {X-NUCLEO-53L8A1};
\node[text=white, font=\tiny, align=center] at (-0.45,0.4)
  {VL53L8CX, 8 x 8 zones\\volatile program memory};
\node[draw=white, fill=black!60, minimum width=5.5mm, minimum height=5.5mm] (tof) at (1.45,0.35) {};
\node[text=white, font=\tiny] at (1.45,0.35) {ToF};

\node[font=\tiny, align=right, anchor=south east] at (-3.05,0.95) {ST-LINK\\micro USB};
\node[draw=black!70, fill=gray!60, minimum width=4mm, minimum height=6mm] (usb) at (-2.75,0.6) {};

\node[ext, left=17mm of nuc, text width=22mm, minimum height=13mm] (pc)
  {Host PC\\{\tiny build flash decode}};
\draw[usbcable] (usb.west) to[out=180,in=0] (pc.east);

\node[ext, right=19mm of nuc, text width=24mm, minimum height=16mm] (ppk)
  {nRF-PPK2\\{\tiny current plus digital in D0}};
\draw[cable, draw=orange!70!black] (2.6,0.15) -- node[above, font=\tiny]{marker} (ppk.west);

% the open question, drawn as a callout rather than asserted
\node[umlnote, text width=42mm, anchor=north west] at (-4.4,-1.8)
  {Bus selection: jumper or solder bridge? Read UM3120 and look at the board.
   There is no soldering iron on this bench, so a bridge means this chapter is
   I2C only and the faster bus is described rather than built.};
\node[umlnote, text width=42mm, anchor=north west] at (0.6,-1.8)
  {Nothing else is placed by hand. The bus, the data-ready line and the reset
   line all reach the microcontroller through the header; which pins they reach
   is read from UM3120 rather than assumed.};
\end{tikzpicture}
